\section{Quantitative transversality and the absence of a bandwidth-only bound}
\label{sec:transversality}
This section consolidates the previously recorded one-mode transversality
argument. It treats arbitrary smooth backgrounds and does not classify every
zero of the second resonance. The last proposition explains why uniform
polynomial coercivity on larger Fourier boxes is a different, false target.

Let $A\in\mathrm{SL}(2,\mathbb Z)$ be hyperbolic and let g be real, even,
smooth and of mean zero. For $h_g=\{g,\mathcal C_Ag\}/(4\pi)$ let
\[
 Q(g)=\sqrt2\biggl(\sum_{w\in\mathcal O/\{\pm1\}}a_w(h_g)\biggr)_{\mathcal O},
 \qquad \Gamma_A(g)=\|Q(g)\|_{\ell^2}^2.
\]
The weighted Fourier absolute sums make Q a continuous homogeneous quadratic
map. For a nonzero integer mode $b=b_v$, put
\[
 \gamma_b=\Gamma_A(b)=4\pi^2\det(v,Av)^2>0.
\]
The two generated frequencies $v-Av,v+Av$ have opposite nonzero values of
the invariant form $q_A(w)=\det(w,Aw)$; hence their orbit sums cannot cancel.
A rational eigenline of A would have an integral eigenvalue with integral
inverse, necessarily $\pm1$, so $\det(v,Av)\ne0$.

\begin{theorem}[Uniform one-mode sublevel estimate]
For every such g and every $\epsilon>0$,
\[
 \operatorname{Leb}\{t\in\mathbb R:\Gamma_A(g+tb)\leq\epsilon\gamma_b\}
       \leq 2\sqrt2\,\epsilon^{1/4}.
\]
The constant is independent of the background's number of Fourier modes,
coefficients and frequency radius. It is a length estimate in the explicitly
normalized scalar coefficient t, not a lower bound at a prescribed t.
\end{theorem}
\begin{proof}
Project $Q(g+tb)$ onto $Q(b)$. The normalized projection is a monic real
quadratic $p(t)=t^2+\beta t+\eta$, and Cauchy--Schwarz gives
$\Gamma_A(g+tb)\geq\gamma_b p(t)^2$. Completing the square reduces
$|p(t)|\leq\delta$ to $|x^2-a|\leq\delta$. If $a\leq\delta$ this is empty
or an interval of length at most $2\sqrt{2\delta}$. If $a>\delta$, the total
length is
$2(\sqrt{a+\delta}-\sqrt{a-\delta})\leq2\sqrt{2\delta}$.
Take $\delta=\sqrt\epsilon$.
\end{proof}

If T has a density bounded by M, the corresponding probability is at most
$\min(1,2\sqrt2 M\epsilon^{1/4})$. This also holds conditionally on an
arbitrary random background whenever the selected coefficient's conditional
density is bounded by M. No independence of all the other coefficients is
required. For $Z=\Gamma_A(g+Tb)/\gamma_b$, integrating the sublevel estimate
at its transition point $(2\sqrt2 M)^{-4}$ gives
\[
 \mathbb E Z^{-p}\leq\frac{(2\sqrt2 M)^{4p}}{1-4p},
                 \qquad 0<p<\tfrac14.
\]
The exponent threshold is exact for a background-uniform statement: at g=0,
$Z=T^4$, so the negative moment diverges for $p\geq1/4$ when T has a density
bounded below near zero. This statement is about a classical variance
functional, not eigenvalue statistics of a randomly perturbed propagator.

\begin{proposition}[A deterministic three-point certificate]
For every r>0,
\[
 \max\{\Gamma_A(g-rb),\Gamma_A(g),\Gamma_A(g+rb)\}
                    \geq \tfrac14r^4\gamma_b.
\]
If $\Gamma_A(g)=0$, then
$\max\{\Gamma_A(g-rb),\Gamma_A(g+rb)\}\geq r^4\gamma_b$.
\end{proposition}
\begin{proof}
For the monic quadratic above, $p(r)+p(-r)-2p(0)=2r^2$. At least one of its
three absolute values is at least $r^2/2$. For the stronger assertion use
$Q(g+rb)+Q(g-rb)=2r^2Q(b)$ when Q(g)=0 and the triangle inequality.
\end{proof}

\begin{proposition}[No polynomial-bandwidth coercivity on all backgrounds]
For the Pell matrix T, let $\mathcal P_R$ be the real even mean-zero Fourier
polynomials with $\|v\|_\infty\leq R$. There are $g_R\in\mathcal P_R$ with
norm bounded below away from zero and bounded in every fixed smooth seminorm,
such that
\[
             0<\Gamma_T(g_R)=O_L(R^{-L})\quad\hbox{for every }L>0.
\]
Thus no constants c,p>0 give $\Gamma_T(g)\geq cR^{-p}\|g\|_2^4$ for all
nondegenerate g in all these Fourier boxes.
\end{proposition}
\begin{proof}
Choose a nonzero mean-zero smooth generator g of the separated-support
construction, so $\{g,\mathcal C_Tg\}=0$. Let $P_Rg$ be its rectangular
Fourier partial sum and $e_R=P_Rg-g$. Smoothness implies
$M_1(e_R)=O_L(R^{-L})$ for every L and $M_1(P_Rg)\leq M_1(g)$.
The exact Fourier bracket bound
$\|\widehat{\{a,b\}/(4\pi)}\|_1\leq\pi M_1(a)M_1(b)$ gives
\[
 \left\|\widehat{\{P_Rg,\mathcal C_TP_Rg\}/(4\pi)}\right\|_1
       \leq2\pi\|T\|_{\mathrm{op}}M_1(g)M_1(e_R).
\]
For a real even h, $\sigma_T^2(h)\leq\tfrac12\|\widehat h\|_1^2$, so
$\Gamma_T(P_Rg)$ is superpolynomially small. Partial sums converge in every
smooth seminorm and in $L^2$.

To ensure strict positivity, choose one nonzero mode v of radius one.
The polynomial $z\mapsto\Gamma_T(P_Rg+zb_v)$ has degree four and positive
leading coefficient $\Gamma_T(b_v)$. At least one of the five numbers
$e^{-R},2e^{-R},3e^{-R},4e^{-R},5e^{-R}$ is not a root. Choose it as $z_R$.
Then $g_R=P_Rg+z_Rb_v$ has strictly positive Gamma; this exponentially small
change preserves every displayed estimate. End of proof.

The certificate for the fixed radius-eight space remains consistent with this
proposition: its positive coercivity constant depends on that fixed radius.
Neither result settles exact nondegeneracy in every finite Fourier box.
